% !TeX root = ../main.tex
% Zeile 1 legt fest das immer die main.tex kompiliert wird

% ===============================================================
% KOMPLIERHINWEISE FÜR GLOSSAR UND ABBKÜRZUNGSVERZEICHNIS (VS CODE)
% ===============================================================

% Damit das Abkürzungsverzeichnis mit dem Paket glossaries korrekt erzeugt wird,
% reicht ein normaler pdflatex-Durchlauf NICHT aus. Du musst makeglossaries mit ausführen.

% Schritt-für-Schritt – Projektbezogene Lösung (empfohlen):

% 1. Erstelle im Projektverzeichnis einen Ordner ".vscode"
% 2. Lege darin eine Datei "settings.json" an
% 3. Füge dort diesen Inhalt ein:

% {
%   "latex-workshop.latex.recipes": [
%     {
%       "name": "pdflatex -> makeglossaries -> pdflatex*2",
%       "tools": [
%         "pdflatex",
%         "makeglossaries",
%         "pdflatex",
%         "pdflatex"
%       ]
%     }
%   ],
%   "latex-workshop.latex.tools": [
%     {
%       "name": "pdflatex",
%       "command": "pdflatex",
%       "args": [
%         "-interaction=nonstopmode",
%         "-file-line-error",
%         "%DOC%"
%       ]
%     },
%     {
%       "name": "makeglossaries",
%       "command": "makeglossaries",
%       "args": ["%DOCFILE%"]
%     }
%   ]
% }

% 4. Speichere die Datei und starte VS Code neu
% 5. Drücke Strg + Shift + P → "LaTeX Workshop: Build with recipe"
% 6. Wähle das Rezept: "pdflatex -> makeglossaries -> pdflatex*2"

% Ergebnis:
% - Das Abkürzungsverzeichnis wird korrekt erstellt
% - \printglossary zeigt die Einträge aus abkuerzungen.tex an